\chapter{User research and human-centred design}
\label{vol11:ch03}

\epigraph{The user the design serves is the actual user; the user the design assumes is the user the design serves only if the assumption was checked.}{}

\chapmeta{Half-life: medium-life. Archetypes: interface, ethics, uncertainty. Exercise target: 18. Project track: hybrid.}

\noindent
User research is the discipline by which the wants that enter the requirements document are verified rather than assumed. Designs that proceed from assumed wants serve the design organisation's mental model of the user; designs that proceed from verified wants serve the actual user. The two are different; the gap between them is the substantive content of the user-research discipline. This chapter develops what user research is, how it is conducted, and how its findings become engineering input.

\section{What user research is and what it is not}\pathtag{core}

User research is the systematic investigation of who uses (or will use) a design, what they need, how they actually behave, and what the gap between their stated and actual needs is. The investigation is empirical (it studies real users in real contexts) rather than speculative (the designer imagines what users want). The output is evidence that informs design decisions.

User research is not:

\textbf{Marketing research.} Marketing research asks "will users buy this?" or "how should we position this product?"; the questions are about the market response to an offering. User research asks "what do users need?" and "how do users actually use this?"; the questions are about the underlying need and behaviour. The two activities overlap and sometimes share methods (surveys, focus groups) but their primary questions differ.

\textbf{Validation of a design.} Validation of a design (Chapter 2) asks whether the completed design serves the user need. User research at the start of a project asks what the user need is. The two are complementary: user research informs the requirements; validation tests the design against the requirements; user research and validation share methods but at different stages and with different purposes.

\textbf{Designer intuition.} The designer's intuition about what users want is sometimes correct; it is sometimes not. User research is the discipline that tests the intuition against evidence. Designers who skip the research and rely on intuition produce designs that match their own mental model; the mental model is correct for some users (typically users similar to the designer) and incorrect for others.

\textbf{Token engagement.} Some design organisations conduct token user research: a small number of interviews; a survey with leading questions; a focus group that the design team observes but does not engage with. The activity produces documents and slides but does not produce evidence that changes design decisions. The token engagement is worse than no research: it consumes effort and creates the appearance of evidence-based design without the substance.

The substantive form of user research has several characteristics:

\textbf{Conducted by trained researchers.} User research is a discipline; doing it well requires training in interview technique, in observation, in analysis, in inference from small samples. Engineers without the training can do user research, but the quality is variable; design organisations that take user research seriously employ trained researchers or train their engineers.

\textbf{Conducted in context.} The research observes users in the contexts in which they actually use (or will use) the design: their workplace, their kitchen, their car, their hospital ward, their factory floor. Context surfaces requirements that lab settings miss; lab studies systematically underestimate the variability of real-world conditions.

\textbf{Influential on design decisions.} The research findings influence what gets built. Findings that are filed away without affecting design are not user research's product; they are research overhead. The integration of research findings into design decisions is the substantive deliverable.

\textbf{Documented and shareable.} The findings are documented in a form the design team can use: insight summaries, journey maps, personas, requirement implications. The documentation makes the findings persistent and shareable beyond the immediate research engagement.

\begin{principle}[Research the actual user]
The design serves the actual user, not the imagined user. The discipline of user research is the discipline of replacing the imagined user with evidence about the actual user. Designs that proceed from imagined users may succeed (when the imagination is accurate) but they succeed by accident; designs that proceed from researched users succeed by design.
\end{principle}

\section{Methods: interviews, contextual inquiry, ethnography}\pathtag{core}

The qualitative user-research methods produce rich, contextual evidence at modest sample sizes. The methods:

\textbf{User interviews.} The researcher conducts structured or semi-structured conversations with individual users. The interview is typically 30 to 90 minutes; the structure depends on the research questions; the recording (audio or video, with consent) preserves the evidence for analysis. The interview technique includes open-ended questioning (avoiding yes/no questions and leading questions), active listening (following the participant's framing rather than imposing the researcher's), and probing (asking "tell me more about that" or "can you give me an example" to surface the underlying experience). The output is interview notes and transcripts that subsequent analysis distils.

\textbf{Contextual inquiry.} The researcher observes the user performing the work the design will serve, in the user's actual environment, while engaging in a "master-apprentice" dialogue (the user is the master of their work; the researcher is the apprentice learning the work). The technique was formalised by Holtzblatt and Beyer; it is a hybrid of interview and observation that surfaces requirements neither alone would. The output is field notes documenting both the observed behaviour and the user's articulation of what they are doing and why.

\textbf{Ethnography.} The researcher embeds in the user's environment over an extended period (days to months), observing the work, the social structure, the workarounds, the constraints, the values that shape how the work is done. The technique is borrowed from anthropology; in engineering contexts, it produces deep contextual understanding that other methods cannot reach. The cost is the time investment and the requirement for trained ethnographers; the value is the understanding of why things work the way they do.

\textbf{Diary studies.} Users record their own experience over an extended period (typically 1-4 weeks), capturing instances of use, problems encountered, workarounds developed, contexts of use. The diary entries can be text, photos, video, voice memos. The technique captures longitudinal patterns that one-time observation misses; it depends on user discipline and the simplicity of the recording method.

\textbf{Focus groups.} The researcher convenes a group of users (typically 6-10) for a structured discussion. The format is efficient (multiple users in one session) and can surface group dynamics (consensus, disagreement, collective concerns); it is unreliable for individual user behaviour because of social-pressure effects. Focus groups work better for marketing research than for user research.

\textbf{Card sorting.} Users sort cards representing concepts or functions into groups that make sense to them. The technique produces evidence about users' mental models that other methods do not surface; it is particularly useful for information-architecture decisions (how should the menu structure be organised; what categories does the user expect).

\textbf{Co-design workshops.} Users and designers work together to sketch, prototype, or critique design candidates. The technique surfaces user requirements through the user's active engagement with concrete design proposals; it works well when users have substantial domain expertise (clinical users for medical devices; pilots for cockpit interfaces).

The qualitative methods produce small-sample evidence (typically 5 to 30 participants) that is rich and contextual. The sample-size mathematics differs from quantitative research: Nielsen's classic argument is that 5 users surface 80\% of the usability problems with a given task; the diminishing-returns curve flattens after 15-20 users for most studies. The exception is when the user population is heterogeneous (different segments have very different needs); then each segment requires its own sample. The design's research budget is allocated to cover the substantive segments.

\section{Surveys, A/B testing, analytics, telemetry}\pathtag{standard}

The quantitative user-research methods produce statistical evidence at larger sample sizes. The methods:

\textbf{Surveys.} The researcher administers a structured questionnaire to a large sample. The survey can measure user characteristics, preferences, behaviours, satisfaction. The technique is efficient at scale but has known limitations: stated preferences differ from revealed preferences; survey responses are subject to social-desirability bias and to question-framing effects; non-response is systematic (some user types respond more than others). The discipline of survey design includes question wording, response scaling, pilot testing, and statistical analysis of the results.

\textbf{A/B testing.} The researcher deploys two (or more) variants of a design feature to different user groups and measures the difference in outcomes. The technique is the gold standard for quantitative user research when applicable: it produces causal evidence about the variant's effect under realistic conditions. The applicability requires that the variants be deployable (typically easier for software than for hardware), that the outcome metric be measurable, and that the user populations be comparable. The discipline includes power analysis (how large a sample is needed to detect the effect), variance reduction techniques, and care about the interpretation of statistically-significant but practically-trivial effects.

\textbf{Analytics.} The researcher analyses logs of user behaviour with the existing system: which features are used, how often, by whom, in what sequences. The technique produces evidence about what users actually do (revealed preference) rather than what they say they do (stated preference). The technique has substantial coverage (all users; all interactions) but limited context (the analytics show what users did but not why); it is most valuable in combination with qualitative methods that explain the patterns.

\textbf{Telemetry.} The researcher instruments the system (the design under study or the predecessor) to record specific behaviours of interest. Telemetry differs from generic analytics in being purpose-built for the research question; it captures the specific signals the analysis requires rather than the broad behavioural traces analytics captures.

\textbf{Heatmaps and session recordings.} The researcher captures screen recordings of user sessions (with consent) and aggregates them into visualisations (heatmaps of click locations; flow diagrams of navigation paths). The technique surfaces interaction patterns that aggregate statistics miss.

\textbf{Benchmarks.} The researcher measures user performance against established baselines (time on task, error rates, completion rates). The technique produces evidence about how the design compares to alternatives; it is particularly useful for accessibility and usability evaluation.

The quantitative methods produce large-sample evidence (typically hundreds to millions of users) that is statistical and (with care) causal. The methods have substantial limitations: they measure what is measurable; they do not measure what users could not articulate or what the instrumentation did not capture. The combination with qualitative methods is necessary for interpretation; statistics without context can mislead.

\section{Personas, journey maps, jobs-to-be-done}\pathtag{standard}

The user-research findings are synthesised into design-relevant artefacts:

\textbf{Personas.} A persona is a fictional but evidence-grounded representation of a user segment. Each persona has: a name and a photograph (giving the design team something to hold in mind); demographic and contextual characteristics; goals, motivations, frustrations relevant to the design; specific scenarios in which they use the design. Personas are designed to be referred to throughout the design process ("how would Maria use this feature?"); they make the user concrete and present in the design team's deliberations.

The discipline of personas includes: grounding each persona in research evidence rather than designer imagination; representing the population's variation (multiple personas covering different segments); avoiding stereotypes that mask important variation within segments; updating the personas as research evidence accumulates.

\textbf{Journey maps.} A journey map is a visual representation of the user's experience over time, often a process the design must support (e.g., a hospital patient's journey from admission to discharge; a customer's journey from product discovery to purchase to use to disposal). The map identifies: stages of the journey; actions the user takes at each stage; emotions and mindset at each stage; pain points and opportunities at each stage; touchpoints with the design.

Journey maps surface design implications that point-in-time analysis misses: the design that excels at one stage may underperform at others; the design's integration across stages may be the substantive user experience. Journey maps are particularly useful for service design and for complex multi-touchpoint products.

\textbf{Jobs to be done (JTBD).} The JTBD framework, popularised by Christensen, frames the user's needs as "jobs" the user is trying to get done: "I want to ensure my children stay healthy" (the job); the product (the children's vitamins) is the user's means to the job. The framing focuses design attention on the underlying job rather than the immediate product. JTBD analysis produces job statements that requirements are derived from; it complements personas (who is doing the job) and journey maps (how the job unfolds).

\textbf{Empathy maps.} An empathy map captures what a user (or persona) says, thinks, does, and feels relative to the design domain. The map is a quick synthesis technique that makes the user's perspective concrete and shareable; it works well for initial alignment of a design team.

\textbf{Mental models.} A mental-model diagram captures the user's conceptual model of the domain: how they think about the problem the design addresses; what concepts and relationships they understand. The discipline includes contrasting the user's mental model with the system model the design implements; gaps between the two are sources of confusion and error.

The synthesis artefacts are means, not ends: their value is in the design decisions they shape, not in their existence. Design organisations that produce elaborate persona documents that no one consults have completed a research-overhead activity; design organisations that produce minimal but consulted personas have done substantive synthesis.

\section{Accessibility and inclusive design as research practice}\pathtag{core}

Accessibility is the practice of designing so that users with disabilities can use the design effectively. Inclusive design is the broader practice of designing so that the design serves the widest plausible range of users, including those who have been historically excluded from the design's user-research samples.

The two are research practices because they require studying users whose needs the dominant user population may not surface. The user with low vision uses interfaces differently from the user with typical vision; the user with limited dexterity uses controls differently from the user with full dexterity; the user with cognitive variation processes information differently from the user with neurotypical processing; the user who is not native to the design's language community uses the interface differently from the native speaker; the user in a low-resource context uses the design differently from the user in a high-resource context.

The methods include:

\textbf{Research with disability communities.} The user-research samples include people with the relevant disabilities; the research methods are adapted to the participants' communication and interaction needs; the findings inform design decisions that affect those users.

\textbf{Accessibility audits.} The completed design is evaluated against accessibility standards (WCAG for web; Section 508 for US federal procurement; the various national standards) and against the actual experience of users with disabilities.

\textbf{Persona-spectrum analysis.} For each impairment relevant to the design (visual, motor, cognitive, hearing, speech), the design team considers the spectrum of user needs (permanent disability; temporary disability; situational limitation). The framework recognises that designs that work for permanent disability typically also work for temporary disability (a broken arm) and for situational limitation (carrying a child while operating a device).

\textbf{Assistive-technology testing.} The design is tested with the assistive technologies users with disabilities employ: screen readers; alternative keyboards; switch controls; magnification; voice control. The test reveals interaction patterns that ordinary testing misses.

\textbf{Inclusive-design panels.} A panel of users representing the design's relevant population diversity reviews the design at key points; the panel provides ongoing input rather than one-time evaluation.

The accessibility-as-research framing places accessibility within the design discipline rather than as an end-of-process add-on. The end-of-process treatment (the design is built, then evaluated for accessibility, then patched for compliance) reliably produces designs that are inferior to designs that were built with accessibility considered from the start. The discipline of inclusive research at the start is the discipline that produces designs serving the actual range of users rather than the narrow range the designer assumed.

The settled editorial decision for this volume's treatment (Q47 and surrounding decisions) is that accessibility is engineering practice rather than a compliance exercise; the present chapter integrates it accordingly. Chapter 10 returns to accessibility as a lifecycle consideration alongside manufacturability, maintainability, and sustainability.

\begin{archetype}[interface]
User research is the interface archetype applied to the design-research relationship. The interface between the design organisation and the user population determines what gets known about user needs and what gets ignored. Research methods that produce rich interface (qualitative; contextual; with users not yet represented in the population) produce more accurate engineering input than methods that produce thin interface (token engagement; survey-only; representative-of-a-narrow-segment).
\end{archetype}

\section{Synthesis: insights to requirements}\pathtag{standard}

User research produces evidence; the evidence must be synthesised into engineering input. The synthesis is a distinct activity from the research itself.

The synthesis steps:

\textbf{Affinity analysis.} The researcher (typically with the design team) clusters the research findings (interview quotes; observation notes; survey responses; analytics patterns) into themes. The clustering is bottom-up: similar findings group together; groups are labelled with the underlying theme. The technique surfaces patterns across the data that no single observation would reveal.

\textbf{Insight extraction.} For each theme, the researcher articulates an insight: a statement of what the research evidence implies about user needs or behaviour, that is specific enough to act on. "Users find the current process frustrating" is a finding; "users abandon the registration process at the password-strength step, citing the multiple-criteria requirements as opaque" is an insight. Insights are the units of synthesis that feed design decisions.

\textbf{Implications mapping.} Each insight is mapped to its design implications: what the insight suggests should change in the design; what alternative concepts the insight enables or rules out; what additional research the insight motivates. The mapping is the bridge from research to design.

\textbf{Requirement derivation.} Some insights map directly to requirements: the insight that users cannot find the search function leads to a requirement that the search be prominently visible from the home screen. The requirement language is the language of Chapter 2; the insight is the engineering rationale.

\textbf{Hypothesis generation.} Some insights generate hypotheses that subsequent research or testing can confirm or refute. The hypotheses become the agenda for follow-up research and for design iteration.

\textbf{Communication artefacts.} The insights and implications are documented in forms the design team uses: insight summaries (1-page documents per insight); journey maps; personas; design principles derived from cross-cutting insights. The artefacts make the research persistent and shareable beyond the immediate engagement.

The synthesis is where user research's value is realised. Research that produces evidence but is not synthesised produces overhead; synthesis that produces design implications that are not acted on produces theatre. The integration of synthesised research into design decisions is the substantive deliverable.

\section{Failure: research that asked the wrong questions; the wrong-user trap}\pathtag{core}

User research can be done badly. The failure modes:

\textbf{Wrong questions.} The research investigates questions that do not bear on substantive design decisions. The questions may be too broad ("what do users want?"), too narrow ("which colour should the button be?"), or off-topic for the design at hand ("how do users feel about our brand?"). The research produces evidence but not engineering input. The discipline of research planning is the discipline of identifying the substantive questions before the research begins.

\textbf{Wrong-user trap.} The research samples users who are convenient or who are similar to the designer rather than users who represent the design's actual population. The convenience sample produces evidence about itself but not about the actual user population. The trap is most consequential when the convenience sample systematically excludes the users with the greatest needs: the accessibility considerations are not surfaced because no users with disabilities were in the sample; the international considerations are not surfaced because no users from outside the design organisation's home market were sampled; the low-income considerations are not surfaced because no users in low-income contexts participated.

\textbf{Leading questions.} The interview questions or survey items signal the answer the researcher expects; the participants provide the expected answer. The research produces evidence that confirms what the researcher already believed; the design decisions follow the confirmation rather than substantive findings. The discipline of question design (open-ended; neutral phrasing; explicit alternatives) is the discipline that resists the leading-questions failure.

\textbf{Selection-on-the-dependent-variable.} The research sample is selected based on the outcome being studied: only users who use the product are interviewed about what users want from the product; only patients who complete treatment are surveyed about treatment quality. The selection biases the findings; the users who do not use the product (the ones whose absence the design might fix) are not heard. The discipline of sampling includes the discipline of sampling outside the current user base when the design's success depends on attracting users beyond it.

\textbf{Confirmation-driven analysis.} The research data are analysed in ways that confirm the design team's hypotheses; data that contradict the hypotheses are dismissed as outliers or as misinterpretations. The bias is well-documented in human cognition; the discipline of structured analysis (with explicit alternative hypotheses; with independent analyst review; with quantitative tracking of supporting versus refuting evidence) is the institutional response.

\textbf{Insight inflation.} A small number of observations are presented as if they were established findings; the design decisions are based on the over-generalised insight; the actual user population does not behave as the insight predicted. The discipline of communicating evidence with explicit confidence bounds is the response; insights should be qualified by the sample size, the variation observed, and the conditions under which the evidence was collected.

\textbf{Token research.} The research is conducted in form but not in substance; the findings are filed but do not influence design decisions; the user-research line item is checked off the project plan without producing engineering input. The remediation is institutional commitment to research as a substantive input rather than a procedural requirement.

The Therac-25 case (Volume X Chapter 5) is partly a user-research failure: the operator interface that allowed the fast-data-entry sequence to expose the race condition was not validated against actual operators' typing patterns; the operators learned to type quickly to keep up with the workflow; the design's assumption that operators would type slowly enough to avoid the race was not checked against the actual user behaviour.

The Boeing 737 MAX MCAS case (Volume X Chapter 8, Chapter 9) is partly a user-research failure: the pilot population's actual response to MCAS activation under realistic stress was not characterised; the design assumed pilots would recognise the MCAS behaviour as runaway-stabiliser and apply the memory-item procedure; the actual response under the conditions of the Lion Air and Ethiopian flights diverged.

\section{Worked examples: medical-device user research, industrial-control-room redesign}\pathtag{standard}

\subsection*{Medical-device user research}

A new infusion pump is being designed for use in intensive-care units. The user population includes registered nurses, physicians, pharmacists, biomedical technicians, and patients (whose interactions are limited but include critical patient-controlled-analgesia use). The use environment is the ICU bedside: noisy, interrupted, multi-patient assignments, sleep-deprived shifts, high-stakes outcomes, integration with the hospital information system, regulated by FDA.

The user research program might include:

\textbf{Contextual inquiry in ICUs.} Researchers observe nurses programming current-generation infusion pumps in actual bedside use. The observations surface: the actual typing patterns under workload; the actual interruption rates; the workarounds nurses use when the interface does not match their workflow; the error patterns and recovery patterns. The output is field-notes documenting the observed practice and insight summaries.

\textbf{Interviews with multiple user roles.} Nurses (the primary users); physicians (the prescribers); pharmacists (the medication-preparation experts); biomedical technicians (the maintainers); patients with patient-controlled analgesia (the end recipients). Each role surfaces requirements the other roles miss.

\textbf{Usability testing with the current pump.} Representative nurses perform standard tasks with the current pump under controlled conditions; the testing surfaces specific interaction problems and their consequences.

\textbf{Heuristic evaluation by accessibility experts.} The current pump is evaluated against accessibility standards and against the needs of nurses with motor or visual variation.

\textbf{Analysis of incident reports.} The FDA MAUDE database and the hospital's internal incident-reporting system are mined for events involving infusion-pump misuse; the patterns inform requirements for the new design.

The synthesis produces: personas for the major user roles; journey maps for the medication-administration process; insights about specific failure modes; requirements that address the insights; design principles for the next-generation pump.

The requirements that emerge might include: programming-interface design that resists fast-typing race conditions; explicit confirmation of dose and rate against patient parameters; integration with the hospital pharmacy system to verify orders; visual and audible status that nurses can assess from across the room; logging that supports post-event analysis; failure modes that default to safe states.

\subsection*{Industrial-control-room redesign}

An ageing chemical-plant control room is being redesigned to address operator workload, alarm flood, situational awareness, and integration with new instrumentation. The user population is the operating crews (12-hour shifts; rotating personnel; mixed experience). The use environment is the central control room: continuous operation; high-consequence decisions; coordination across multiple operators.

The user research program might include:

\textbf{Ethnography in the existing control room.} Researchers observe operators over extended shifts: the steady-state monitoring; the upset response; the shift handover; the abnormal-condition diagnosis. The ethnography surfaces the social structure of operations, the communication patterns, the use of the existing displays, and the workarounds that operators have developed.

\textbf{Cognitive task analysis.} The researchers decompose the operators' diagnostic tasks into the cognitive operations they perform: pattern recognition; differential diagnosis; verification; intervention selection; outcome monitoring. The analysis informs display design: which information must be visible at a glance; which must be accessible on demand; which must trigger active alerts.

\textbf{Simulator-based observation.} The operators perform standard upset scenarios in a high-fidelity simulator; the observation surfaces the actual diagnostic and intervention patterns under realistic conditions.

\textbf{Alarm-flood analysis.} The current alarm system is analysed for false-positive rates, prioritisation, and operator response patterns. The analysis informs alarm-management requirements (consolidation, prioritisation, suppression of nuisance alarms).

\textbf{Crew-resource-management observation.} The team coordination during normal and abnormal operations is observed; the observations inform requirements for inter-operator communication, role definition, and procedural support.

The synthesis produces: insights about the gap between the designed work and the actual work; requirements for the new control room's display, alarm, and procedural support; design principles for the integration with the new instrumentation.

The requirements that emerge address the lessons from Volume IX (control), Volume X Chapter 9 (human factors), and Volume X Chapter 13 (weak-signal recognition): displays that support skill-based, rule-based, and knowledge-based operation; alarm prioritisation that resists flood; cross-operator situational-awareness support; procedural support that does not constrain operator judgment in novel situations.

\begin{principle}[Research before design; verify after]
User research at the start of the design replaces designer assumptions with evidence about the user population. Validation at the end of the design confirms that the executed design serves the population the research characterised. The two activities together are the discipline by which design produces artefacts that serve actual users rather than imagined ones.
\end{principle}

\begin{project}[User research for one product]
Select one product or service the reader uses regularly. Conduct a small user-research study: interview 3-5 users (or, if no other users are accessible, conduct the equivalent of self-ethnography on the reader's own use over a week); observe at least one user (or yourself) in actual use; synthesise the findings into 5 specific insights, each grounded in evidence from the research; for each insight, identify a design implication (what would change in a redesign). Deliverable: a 5-page document with the research description, the insights with evidence, and the implications.

\textbf{Hazard.} The exercise's value depends on the discipline of the research methods, not on the elaborateness of the deliverable. A small but well-conducted study produces actionable insights; a large but poorly-conducted study produces overhead.
\end{project}

\section{Exercises}

\subsection*{Calculation}\pathtag{core}

\begin{exercise}
Nielsen's classic argument is that 5 users surface approximately 80\% of usability problems in a study with a given task. Using the model that each user surfaces each problem independently with probability $p$, derive $p$ from the 80\%-at-5-users figure and predict the fraction surfaced at 10, 15, and 20 users.
\end{exercise}

\begin{exercise}
An A/B test must detect an effect size of $\Delta$ with power 0.8 and significance 0.05. Compute the required sample size per arm for $\Delta / \sigma = 0.1$, $0.5$, $1.0$. Discuss the practical implications for designing A/B tests that can detect substantive effects.
\end{exercise}

\begin{exercise}
A user-research study budget is 200 person-hours. Each interview takes 4 hours (preparation, conduct, transcription, analysis); each contextual inquiry session takes 8 hours; each ethnographic day takes 12 hours. Compute the maximum number of each type of engagement under the budget, and discuss how to allocate the budget across methods.
\end{exercise}

\subsection*{Derivation}\pathtag{standard}

\begin{exercise}
For a user population stratified into $K$ segments with proportions $p_1, \ldots, p_K$ and within-segment variance $\sigma^2$, derive the optimal allocation of a sample of $N$ across segments to minimise the variance of the population-level estimate. Identify when proportional allocation is optimal and when stratified sampling improves on it.
\end{exercise}

\begin{exercise}
Derive the conditions under which qualitative research (small samples; rich data) and quantitative research (large samples; thin data) produce comparable evidence about a design question. Identify the dimensions on which qualitative dominates and those on which quantitative dominates.
\end{exercise}

\subsection*{Estimation}\pathtag{standard}

\begin{exercise}
Estimate the cost of a user-research program adequate to support a medical-device development (regulated by FDA; 18-month development; user populations include nurses, physicians, and patients). Identify the dominant cost categories.
\end{exercise}

\begin{exercise}
Estimate the fraction of failed products (products withdrawn from market within 2 years of launch) where the failure is attributable to inadequate user research versus inadequate execution. Identify the empirical evidence the estimate rests on.
\end{exercise}

\begin{exercise}
Estimate the analytics data volume from a software product with 1 million monthly active users, capturing 50 events per user per session, with 10 sessions per month. Identify the storage, processing, and privacy implications.
\end{exercise}

\subsection*{Simulation}\pathtag{standard}

\begin{exercise}
Implement a simulation of a usability study with $n$ users, each with independent probability $p$ of encountering each of $M$ usability problems. Plot the fraction of problems discovered versus $n$ for $p = 0.31$ (Nielsen's classic figure). Compare to alternative discovery models (e.g., problems with varying severity having different per-user probabilities).
\end{exercise}

\begin{exercise}
Simulate the wrong-user trap: a design is evaluated by a sample of users systematically different from the actual user population. Sweep the systematic difference (e.g., the sample's average is shifted by $\delta$ on the relevant dimension) and identify the conditions under which the evaluation correctly predicts population behaviour.
\end{exercise}

\subsection*{Design}\pathtag{standard}

\begin{exercise}
Design a user-research plan for a hypothetical product (e.g., a smart-home device for elderly users living independently). Specify: research questions; user segments; methods; sample sizes; timeline; synthesis approach; integration with the design process.
\end{exercise}

\begin{exercise}
Design an interview-protocol template for a user-research study. Include: opening rapport-building; warm-up questions; main topic questions with probes; closing questions; informed-consent language; recording and confidentiality provisions.
\end{exercise}

\begin{exercise}
Design an accessibility-evaluation protocol for a software product. Specify: the standards used as reference; the user representatives included; the assistive technologies tested; the synthesis of findings into design implications.
\end{exercise}

\subsection*{Judgement}\pathtag{mastery}

\begin{exercise}
Argue: token user research is worse than no user research. Identify the conditions under which the argument holds and the conditions under which token research is preferable to no research at all.
\end{exercise}

\begin{exercise}
Discuss the wrong-user trap as a structural failure mode of design organisations. Identify the institutional conditions that produce the trap and the institutional changes that would reduce its incidence.
\end{exercise}

\begin{exercise}
Argue: accessibility considerations should be integrated into user research from the start rather than addressed as a post-design compliance exercise. Identify the costs and benefits of each approach and discuss what determines which is appropriate for a given project.
\end{exercise}

\begin{exercise}
A user-research finding contradicts the senior product manager's strongly-held intuition about the user population. Discuss the engineer's options for handling the situation, the professional-ethics implications, and the institutional structures that protect against the suppression of inconvenient research findings.
\end{exercise}

\begin{exercise}
The Therac-25 case (Volume X Chapter 5) had operator-interface assumptions that were not checked against actual operator typing patterns. Discuss what user-research methods would have surfaced the operator's fast-data-entry behaviour in time to affect the design. Identify the institutional barriers that may have prevented this research from being conducted.
\end{exercise}
